%% $Id: pst-func.tex 885 2024-03-31 07:25:14Z herbert $
%%
%% This is file `pst-func.tex',
%%
%% IMPORTANT NOTICE:
%%
%% Package `pst-func.tex'
%%
%% Herbert Voss <hvoss@tug.org>
%%
%% This program can be redistributed and/or modified under the terms
%% of the LaTeX Project Public License Distributed from CTAN archives
%% in directory macros/latex/base/lppl.txt.
%%
%% DESCRIPTION:
%%   `pst-func' is a PSTricks package to plot special functions
%%
%% For a ChangeLog go the the end
%%
\csname PSTfuncLoaded\endcsname
\let\PSTfuncLoaded\endinput
% Requires some PSTricks packages
\ifx\PSTricksLoaded\endinput\else   \input pstricks.tex\fi
\ifx\PSTnodesLoaded\endinput\else   \input pst-plot.tex\fi
\ifx\PSTmathLoaded\endinput \else   \input pst-math.tex\fi
\ifx\PSTtoolsLoaded\endinput \else  \input pst-tools.tex\fi
\ifx\PSTricksAddLoaded\endinput\else\input pstricks-add.tex\fi
\ifx\PSTXKeyLoaded\endinput\else    \input pst-xkey.tex \fi
%
\edef\PstAtCode{\the\catcode`\@} \catcode`\@=11\relax
% interface to the `xkeyval' package
\pst@addfams{pst-func}
%
\def\fileversion{1.02a}
\def\filedate{2024/03/31}
\message{`PST-func' v\fileversion, \filedate\space (hv)}
%
\pstheader{pst-func.pro}
%\pstheader{pst-math.pro}% for GAMMALN
%
% Shortcuts ....

\def\ChebyshevT{ tx@FuncDict begin ChebyshevT end }
\def\ChebyshevU{ tx@FuncDict begin ChebyshevU end }

%
\define@key[psset]{pst-func}{epsilon}[1.e-08]{\def\psk@epsilon{#1 }}
\psset[pst-func]{epsilon=1.e-08}
\define@key[psset]{pst-func}{xShift}[0]{\def\psk@xShift{#1}}
\psset[pst-func]{xShift=0}
%
\define@key[psset]{pst-func}{cosCoeff}[0]{\def\psk@cosCoeff{#1}}
\define@key[psset]{pst-func}{sinCoeff}[1]{\def\psk@sinCoeff{#1}}
\psset[pst-func]{cosCoeff=0,sinCoeff=1} % coeff=a0 a1 a2 a3 ...
%
\def\psFourier{\@ifnextchar[{\psFourier@i}{\psFourier@i[]}}
\def\psFourier@i[#1]#2#3{{%
  \pst@killglue
  \psset{#1}
  \psplot[algebraic=false]{#2}{#3}{%
      /psFtype (cos) def
      /Fourier {
        aload length /n exch def
        n -1 roll 2 div n 1 roll % a0/2
        n 1 sub -1 0 {
          /i exch def
          i x mul 180 mul 3.141592 div
          psFtype (sin) eq {sin}{cos} ifelse
          mul n 1 roll
        } for
        n 1 sub -1 1 { pop add } for
      } def
      [\psk@cosCoeff] Fourier
      /psFtype (sin) def
      [0 \psk@sinCoeff] Fourier add
    }%
}\ignorespaces}
%
\define@key[psset]{pst-func}{coeff}[0 1]{\def\psk@coeff{#1}}
\define@key[psset]{pst-func}{Derivation}[0]{\def\psk@Derivation{#1}}
\define@boolkey[psset]{pst-func}[Pst@]{markZeros}[true]{}
\define@key[psset]{pst-func}{epsZero}[0.1]{\def\psk@epsZero{#1}}
\define@key[psset]{pst-func}{dZero}[0.1]{\def\psk@dZero{#1}}
\define@key[psset]{pst-func}{zeroLineTo}[-1]{\def\psk@zeroLineTo{#1}}
\define@key[psset]{pst-func}{zeroLineColor}[black]{\pst@getcolor{#1}\psk@zeroLineColor}
\newdimen\psk@zeroLineWidth
\define@key[psset]{pst-func}{zeroLineWidth}[0.5\pslinewidth]{\pssetlength\psk@zeroLineWidth{#1}}
\define@key[psset]{pst-func}{zeroLineStyle}[dashed]{%
  \@ifundefined{psls@#1}%
    {\@pstrickserr{Line style `#1' not defined}\@eha}%
    {\edef\psk@zeroLineStyle{#1}}%
}
\psset[pst-func]{%
       coeff=0 1,      % coeff=a0 a1 a2 a3 ...
       Derivation=0, % 0 is the original function
       markZeros=false,% no dots for the zeros
       epsZero=0.1,    % the distance between two zero points
       dZero=0.1,      % the distance of the x value for scanning the function
       zeroLineTo=-1,  % a line to the value of the lineTo's Derivation (-1= none)
       zeroLineStyle=dashed,%
       zeroLineWidth=0.5\pslinewidth,%
       zeroLineColor=black}%
%
\def\psGetZeros{\pst@object{psGetZeros}}
\def\psGetZeros@i(#1,#2)#3{%
  \pst@killglue
  \begin@SpecialObj
  \use@par
  \addto@pscode{
    tx@FuncDict begin
    /xStart #1 def /xEnd #2 def
    /x xStart def
    /func { /x exch def \ifPst@algebraic (#3) AlgParser cvx exec \else #3 \fi } def
    #1 Steffensen
    \ifPst@markZeros \fi
    end
  }%
  \end@SpecialObj
\ignorespaces
}
\iffalse
/Steffensen {% the start value and the function must be on top of the stack
  /y0 exch def % the start value
  /Iter 0 def /MaxIter 200 def
  {
    y0 Function /F exch def
    F abs eps2 lt { exit } if
    y0 F sub /Phi exch def
    Phi Function /F2 exch def
    F2 abs eps2 le { exit }{
      Phi y0 sub dup mul Phi F2 sub 2 Phi mul sub y0 add Div /Diff exch def
      y0 Diff sub /y0 exch def
      Diff abs eps1 le { exit } if
    } ifelse
    /Iter Iter 1 add def
    Iter MaxIter gt { exit } if
  } loop
  y0 % the returned value ist the zero point
  0 \tx@ScreenCoor pstack
  3 0 360 arc gsave 0 0 1 setrgbcolor fill grestore 1 setlinewidth stroke
} def
\fi
%
\def\psPolynomial{\pst@object{psPolynomial}}
\def\psPolynomial@i#1#2{%
  \pst@killglue%
  \begingroup%
  \use@par%
  \@nameuse{beginplot@\psplotstyle}%
  \gdef\psplot@init{}%
  \@nameuse{testqp@\psplotstyle}%
  \addto@pscode{%
    tx@FuncDict begin
    /coeff [ \psk@coeff ] def
    /x0 #1 def /x1 #2 def
    /dx x1 x0 sub \psk@plotpoints\space div def
    /Derivation \psk@Derivation\space def
    /x x0 def
    \ifPst@markZeros
      5 dict begin % hold all local!
      gsave
      \pst@number\psk@zeroLineWidth SLW
      \pst@usecolor\psk@zeroLineColor
      \psk@epsZero\space \psk@dZero\space FindZeros
        aload length {  % zero array is on stack
        /xZero exch def
        xZero \pst@number\psxunit mul /xPixel exch def
        \psk@dotsize
        \@nameuse{psds@\psk@dotstyle}%
        xPixel 0 Dot
        \psk@zeroLineTo\space 0 ge { % line to function \psk@lineTo
          xPixel 0 moveto
          xZero coeff \psk@zeroLineTo\space FuncValue
          \pst@number\psyunit mul xPixel exch L
          \@nameuse{psls@\psk@zeroLineStyle}
        } if
      } repeat
      grestore
      end
    \fi
    /xy {
      x \psk@xShift\space sub coeff Derivation FuncValue \pst@number\psyunit mul
      x \pst@number\psxunit mul exch
    } def
    xy moveto
  }%
  \if@pst%	lines and dots
    \psPolynomial@iii
  \else%	curves
    \psPolynomial@ii
  \fi
  \endgroup
\ignorespaces}
%
\def\psPolynomial@ii{%
  \addto@pscode{%
     xy \@nameuse{beginqp@\psplotstyle}
    \psk@plotpoints {
      xy \@nameuse{doqp@\psplotstyle}
      /x x dx add def
    } repeat
    xy \@nameuse{doqp@\psplotstyle}
    end
  }%
  \@nameuse{endqp@\psplotstyle}%
}
\def\psPolynomial@iii{%    curves
  \addto@pscode{%
    mark
    /n 2 def
    \psk@plotpoints {
      xy
      n 2 roll
      /n n 2 add def
      /x x dx add def
    } repeat
    /x x1 def
    xy
    n 2 roll
    end
  }%
  \@nameuse{endplot@\psplotstyle}%
}
%
% Bessel 2004-06-08
% Manuel Luque, Herbert Voss
% Look at the end for some more documentation about the algorithm
%
\define@key[psset]{pst-func}{constI}[1]{\def\psk@constI{#1 }}
\define@key[psset]{pst-func}{constII}[0]{\def\psk@constII{#1 }}
\psset{constI=1,constII=0}
%
\def\psBessel{\@ifnextchar[{\psBessel@i}{\psBessel@i[]}}
\def\psBessel@i[#1]#2#3#4{{%%%  #2 = n
  \pst@killglue
  \psset{plotpoints=500,constI=1,constII=0}%
  \psset{#1,algebraic=false}%
  \pstVerb{ /toDegrees { 57.29577951 mul } def }
  \parametricplot{#3}{#4}{%
    /J1 0 def
    /xBessel t toDegrees def
    0 0.1 180 {
      /tB exch toDegrees def
      /J1 J1 0.1 xBessel
        tB sin mul tB #2\space mul sub cos mul add def
    } for
    t J1 180 div \psk@constI mul \psk@constII add
  }%
}\ignorespaces}
%
%
\def\psModBessel{\@ifnextchar[{\psModBessel@i}{\psModBessel@i[]}}%% hv 20111021
\def\psModBessel@i[#1]#2#3{{%%%  #2 = n
  \pst@killglue%
  \psset{nue=0,#1}%
  \psplot{#2}{#3}[ /nue \psk@nue def /epsilon 1e-20 def ]{%
    /Sum 0 def
    /Iter 0 def
    {/Sum_Iter
       x dup mul 4 div Iter exp        % nominator
       nue Iter add 1 add GAMMA Iter tx@AddMathFunc begin ! end mul  % denominator
       Div def
     Sum_Iter abs epsilon lt { exit } if
     /Sum Sum Sum_Iter add def
     /Iter Iter 1 add  def
    } loop
    x 0.5 mul nue exp Sum mul
  }%
}\ignorespaces}
%
\define@key[psset]{pst-func}{n}[1]{\def\psk@func@n{#1 }}
\psset[pst-func]{n=0}
%
\def\psLaguerre{\@ifnextchar[\psLaguerre@i{\psLaguerre@i[]}}%% hv 20200111
\def\psLaguerre@i[#1](#2,#3){{%%% n
  \pst@killglue
  \psplot[#1]{#2}{#3}{%
    /L_n 1 def
    0 1 \psk@func@n {
      /k_i ED
      -1 k_i exp k_i factorial div 
      \psk@func@n k_i MoverN mul
      x k_i exp mul L_n add 
      /L_n ED
    } for
    x \pst@number\psxunit mul L_n 
    %dup == %\pst@number\psyunit div
  }%
}\ignorespaces}
%
\def\psLaguerreC{\pst@object{psLaguerreC}}%% hv 20200612
\def\psLaguerreC@i(#1,#2)(#3,#4){{%%% n (xMin,xMax)(yMin,yMax)
  \pst@killglue
  \begin@SpecialObj
  \addto@pscode{%
%    20 dict begin
    /yC #3 def
    /dx #2 #1 sub \psk@plotpoints div def
    /dy #4 #3 sub \psk@plotpoints div def
    /L_i { 
      5 userdict begin
      /X ED
      /n_i ED 
      /L_n 0 def
      0 1 n_i {
        /k_i ED
        -1 k_i exp k_i factorial div 
        n_i k_i MoverN mul
        X k_i exp mul L_n add 
        /L_n ED
      } for
      L_n % return L_n(x)
      end
    } def
    /n \psk@func@n def
    [ % start array of x/y values
    \psk@plotpoints { /xC #1 def \psk@plotpoints { 
    /Re 0 def
    0 1 \psk@func@n 2 div floor {
      /j ED % loop index
      /j2 j dup add def  % 2j
      n j2 sub xC L_i % calculate L_{n-2j}(x)
      j2 exp % L_{n-2j}^{2j}(x)
      -1 j exp yC j2 exp mul % (-1^j)*y^(2j)
      j2 factorial % (2*j)!
      div
      mul
      Re add /Re ED 
    } for
    /Im 0 def
    0 1 \psk@func@n 1 sub 2 div floor {
      /j ED % loop index
      /j2 j dup add def  % 2j
      n j2 sub 1 sub xC L_i % calculate L_{n-2j-1}(x)
      j2 1 add exp % L_{n-2j}^{2j+1}(x)
      -1 j 1 sub exp yC j2 1 add exp mul % (-1^j)*y^(2j+1)
      j2 1 add factorial % (2*j+1)!
      div
      mul
      Im add /Im ED 
    } for
    Re Im 2 copy \tx@ScreenCoor
%    Re Im dup mul exch dup mul add sqrt dup 4 gt {pop 4} if 4 div setgray
%    xC yC \tx@ScreenCoor 0 360 2 arc gsave fill stroke grestore
    xC dx add /xC ED
    } repeat    
      yC dy add /yC ED
    } repeat
    false \tx@NArray 
    \psk@dotsize
    \@nameuse{psds@\psk@dotstyle}
    \tx@setStrokeTransparency
    newpath
    n { transform floor .5 add exch floor .5 add exch itransform Dot  } repeat 
%%    end
  }%
  \end@SpecialObj
}\ignorespaces}
%

\def\psLaguerreCC{\pst@object{psLaguerreCC}}%% hv 20200612
\def\psLaguerreCC@i(#1,#2)(#3,#4){{%%% n (xMin,xMax)(yMin,yMax)
  \pst@killglue
  \begin@SpecialObj
  \addto@pscode{%
%    20 dict begin
gsave
    /yC #3 def
    /dx #2 #1 sub \psk@plotpoints div def
    /dy #4 #3 sub \psk@plotpoints div def
    /n \psk@func@n def
    [ % start array of x/y values
    \psk@plotpoints { /xC #1 def \psk@plotpoints { 
/z [xC yC ] def
/L0 [1 0] def
/L1 [1 z {-1 mul} forall 3 1 roll add exch] def
n 1 le {
 n 0 eq {/L1 L0 def} if
 n 1 eq {/L1 L1 def} if
 }{
/m 1 def
n 1 sub
 {
 /L2 1 m 1 add div
     [2 m mul 1 add z {-1 mul} forall 3 1 roll add exch] L1 cxmul 
      L0 m neg cxrmul cxadd
      exch cxrmul def
 /L0 L1 def
 /L1 L2 def
/m m 1 add def
} repeat
} ifelse
L1 aload pop \tx@ScreenCoor
    xC dx add /xC ED
    } repeat    
      yC dy add /yC ED
    } repeat
    false \tx@NArray 
    \psk@dotsize
    \@nameuse{psds@\psk@dotstyle}
    \tx@setStrokeTransparency
    newpath
    n { transform floor .5 add exch floor .5 add exch itransform Dot  } repeat 
grestore
%%    end
  }%
  \end@SpecialObj
}\ignorespaces}
%

\def\psLaguerreCCC{\pst@object{psLaguerreCCC}}%% hv 20200612
\def\psLaguerreCCC@i(#1,#2)(#3,#4){{%%% n (xMin,xMax)(yMin,yMax)
  \pst@killglue
  \begin@SpecialObj
  \addto@pscode{%
%    20 dict begin
    /yC #3 def
    /dx #2 #1 sub \psk@plotpoints div def
    /dy #4 #3 sub \psk@plotpoints div def
    [ % start array of x/y values
    \psk@plotpoints { 
      /xC #1 def 
      \psk@plotpoints { 
        /Z [xC yC] def
        /L [1 0] def
        0 1 \psk@func@n {
          /k ED
          \psk@func@n k MoverN -1 k exp k factorial div mul 
          Z k 
          tx@FuncDict begin cxexp exch 
          cxrmul L cxadd end /L ED
        } for
        L aload pop \tx@ScreenCoor
        xC dx add /xC ED
      } repeat    
      yC dy add /yC ED
    } repeat
%    ] aload length 3 1 roll moveto 2 div cvi 1 sub { lineto } repeat 1 0 0 setrgbcolor 0.1 setlinewidth stroke 
    false \tx@NArray 
    \psk@dotsize
    \@nameuse{psds@\psk@dotstyle}
    \tx@setStrokeTransparency
    newpath
    n { transform floor .5 add exch floor .5 add exch itransform Dot  } repeat 
%%    end
  }%
  \end@SpecialObj
}\ignorespaces}
%
\def\psLaguerreCCCC{\pst@object{psLaguerreCCCC}}%% hv 20200612
\def\psLaguerreCCCC@i(#1,#2)(#3,#4){{%%% n (xMin,xMax)(yMin,yMax)
  \pst@killglue
  \begin@SpecialObj
  \addto@pscode{%
%    20 dict begin
    /yC #3 def
    /dx #2 #1 sub \psk@plotpoints div def
    /dy #4 #3 sub \psk@plotpoints div def
    % start array of x/y values
    \psk@plotpoints { 
      /xC #1 def 
      [
      \psk@plotpoints { 
        /Z [xC yC] def
        /L [1 0] def
        0 1 \psk@func@n {
          /k ED
          \psk@func@n k MoverN -1 k exp k factorial div mul 
          Z k 
          tx@FuncDict begin cxexp exch 
          cxrmul L cxadd end /L ED
        } for
        L aload pop \tx@ScreenCoor
        xC dx add /xC ED
      } repeat    
      ] 
     aload length 3 1 roll moveto 2 div cvi 1 sub { lineto } repeat 1 0 0 setrgbcolor 0.1 setlinewidth stroke 
      yC dy add /yC ED
    } repeat
%    ] aload length 3 1 roll moveto 2 div cvi 1 sub { lineto } repeat 1 0 0 setrgbcolor 0.1 setlinewidth stroke 
%    false \tx@NArray 
%    \psk@dotsize
%    \@nameuse{psds@\psk@dotstyle}
%    \tx@setStrokeTransparency
%    newpath
%    n { transform floor .5 add exch floor .5 add exch itransform Dot  } repeat 
%%    end
  }%
  \end@SpecialObj
}\ignorespaces}
%
\def\psContourLaguerre{\pst@object{psContourLaguerre}}%% hv 20200612
\def\psContourLaguerre@i(#1,#2)(#3,#4){{%%% n (xMin,xMax)(yMin,yMax)
  \pst@killglue
  \begin@SpecialObj
  \addto@pscode{%
    /yC #3 def
    /dx #2 #1 sub \psk@plotpoints div def
    /dy #4 #3 sub \psk@plotpoints div def
    /contourdata [ % start array of x/y values
    \psk@plotpoints { 
      /xC #1 def 
      [
      \psk@plotpoints { 
        /Z [xC yC] def
        /L [1 0] def
        0 1 \psk@func@n {
          /k ED
          \psk@func@n k MoverN -1 k exp k factorial div mul 
          Z k 
          tx@FuncDict begin cxexp exch 
          cxrmul L cxadd end /L ED
        } for
%        L aload cxnorm exch pop 
        xC yC L cxre exch pop 
        xC dx add /xC ED
      } repeat    
      ]
      yC dy add /yC ED
    } repeat
    ] def     
    /zMax 0 def /zMin 0 def 			% lowest and highest value
    contourdata aload length /N ED % get the no of arrays
    N {						% inside contourdata
      /data ED 					% save first inner array
      data aload length 3 div round cvi {	% get the records
        \psk@colorOffset 
        dup zMin lt 				% z<zMin?
	  { /zMin ED }				% yes, save it
	  { dup zMax gt { /zMax ED }{ pop } ifelse } ifelse
        pop pop					% delete x y 
      } repeat
    } repeat
    clear					% clear stack
    /dz zMax zMin sub def			% value range
    /steps \psk@colSteps\space 0 gt { true }{ false } ifelse def
    0 1 N 1 sub { 				% for i=0 to N-1
      contourdata exch get /data ED		% get first array
      data aload length 3 div round cvi {	
        \psk@colorOffset
        /z ED /y ED /x ED
        z zMin sub dz div 			% relative z (0..1)
        \ifPst@colored 400 mul 380 add tx@addDict begin wavelengthToRGB 
           Red Green Blue end setrgbcolor 	% set color	
        \else
          steps {\psk@colSteps\space mul round \psk@colSteps\space div} if
          setgray 
	\fi
        x \pst@number\psxunit mul y \pst@number\psyunit mul 
        \pst@number\pst@dimb 0 360 arc fill 
      } repeat
    } for
  }%
  \end@SpecialObj
}\ignorespaces}



\def\psLegendre{\@ifnextchar[{\psLegendre@i}{\psLegendre@i[]}}%% hv 20200115
\def\psLegendre@i[#1](#2,#3){{%%% n
  \pst@killglue%
  \psplot[#1]{#2}{#3}{%
    /P_n 0 def
    \psk@func@n 0 eq { /P_1 1 def }{ \psk@func@n 1 eq { /P_1 x def }
       { /P_0 1 def
	 /P_1 x def  
	 2 1 \psk@func@n {
           /k_i ED
           k_i dup add 1 sub x mul P_1 mul  % (2n-1)*x*P_1
           k_i 1 sub P_0 mul sub  % − ( n−1 ) P_{n-2}(x)
           k_i div
           /P_0 P_1 def
           /P_1 ED
         } for 
%n Pn ( x ) = ( 2n − 1 )  P_{n−1}(x) − ( n−1 ) P_{n-2}(x)       
      } ifelse } ifelse
    x \pst@number\psxunit mul P_1 
    %dup == %\pst@number\psyunit div
  }%
}\ignorespaces}
%
\define@key[psset]{pst-func}{sigma}[0.5]{\def\psk@sigma{#1 }}
\define@key[psset]{pst-func}{mue}[0]{\def\psk@mue{#1 }}
\define@key[psset]{pst-func}{nue}[1]{\def\psk@nue{#1 }}
\psset[pst-func]{sigma=0.5,mue=0,nue=1}
%
\def\psGauss{\@ifnextchar[{\psGauss@i}{\psGauss@i[]}}
\def\psGauss@i[#1]#2#3{{%
    \pst@killglue%
    \psset{plotpoints=200}%
    \psset{#1}%
    \psplot[algebraic=false]{#2}{#3}{%
        Euler x \psk@mue sub dup mul 2 div \psk@sigma dup mul div neg exp
	1.0 \psk@sigma div TwoPi sqrt div mul%
    }%
}\ignorespaces}
%
\define@key[psset]{pst-func}{Simpson}[5]{\def\psk@Simpson{#1 }}
\psset[pst-func]{Simpson=5}
%
\def\psGaussI{\pst@object{psGaussI}}
\def\psGaussI@i#1#2{%
  \addbefore@par{plotpoints=200,plotstyle=line}
  \begin@OpenObj%
  \addto@pscode{
    /a #1  def
    /dx #2 #1 sub \psk@plotpoints\space div def
    /b a dx add def
    /scx { \pst@number\psxunit mul } def
    /scy { \pst@number\psyunit mul } def
    tx@FuncDict begin
    /C 1 \psk@sigma div TwoPi sqrt div def
    /SFunc {% x on Stack
      Euler exch \psk@mue\space sub dup mul 2 div \psk@sigma\space dup mul div neg exp C mul
    } def
    end
%    a scx 0 moveto
    a scx 0 \@nameuse{beginqp@\psplotstyle}
    \psk@plotpoints 1 sub {
      a b \psk@Simpson false % a b M useXval on Stack
      tx@FuncDict begin Simpson I end % y value on stack
      scy b scx exch \@nameuse{doqp@\psplotstyle} %lineto
      /b b dx add def
    } repeat
%    stroke
  }%
  \end@OpenObj%
}
%
\def\psSi{\pst@object{psSi}}
\def\psSi@i#1#2{%
  \begin@OpenObj%
  \addto@pscode{
    /x #1  def
    /dx #2 #1 sub \psk@plotpoints\space div def
    /scx { \pst@number\psxunit mul } def
    /scy { \pst@number\psyunit mul } def
    x scx x tx@FuncDict begin Si end scy moveto
    \psk@plotpoints 1 sub {
      x dup scx exch tx@FuncDict begin Si end scy lineto
      /x x dx add def
    } repeat
    stroke
  }%
  \end@OpenObj%
}
\def\pssi{\pst@object{pssi}}
\def\pssi@i#1#2{%
  \begin@OpenObj%
  \addto@pscode{
    /x #1  def
    /dx #2 #1 sub \psk@plotpoints\space div def
    /scx { \pst@number\psxunit mul } def
    /scy { \pst@number\psyunit mul } def
    x scx x tx@FuncDict begin si end scy moveto
    \psk@plotpoints 1 sub {
      x dup scx exch tx@FuncDict begin si end scy lineto
      /x x dx add def
    } repeat
    stroke
  }%
  \end@OpenObj%
}
%
\def\psCi{\pst@object{psCi}}
\def\psCi@i#1#2{%
  \begin@OpenObj%
  \addto@pscode{
    /x #1  def
    /dx #2 #1 sub \psk@plotpoints\space div def
    /scx { \pst@number\psxunit mul } def
    /scy { \pst@number\psyunit mul } def
    x scx x tx@FuncDict begin Ci end scy moveto
    \psk@plotpoints 1 sub {
      x dup scx exch tx@FuncDict begin Ci end scy lineto
      /x x dx add def
    } repeat
    stroke
  }%
  \end@OpenObj%
}
\def\psci{\pst@object{psci}}
\def\psci@i#1#2{%
  \begin@OpenObj%
  \addto@pscode{
    /x #1  def
    /dx #2 #1 sub \psk@plotpoints\space div def
    /scx { \pst@number\psxunit mul } def
    /scy { \pst@number\psyunit mul } def
    x scx x tx@FuncDict begin ci end scy moveto
    \psk@plotpoints 1 sub {
      x dup scx exch tx@FuncDict begin ci end scy lineto
      /x x dx add def
    } repeat
    stroke
  }%
  \end@OpenObj%
}
%
\define@key[psset]{pst-func}{PSfont}[NimbusRomNo9L-Regu]{\def\psk@PSfont{/#1 }}
\define@key[psset]{pst-func}{valuewidth}[20]{\pst@getint{#1}\psk@valuewidth }
\define@key[psset]{pst-func}{fontscale}[10]{\pst@checknum{#1}\psk@fontscale }
\define@key[psset]{pst-func}{decimals}[-1]{\pst@getint{#1}\psk@decimals }
\psset[pst-func]{PSfont=NimbusRomNo9L-Regu,fontscale=10,valuewidth=20,decimals=-1}
%
\def\psPrintValue{\pst@object{psPrintValue}}
\def\psPrintValue@i#1{\expandafter\psPrintValue@ii#1,,\@nil}
\def\psPrintValue@ii#1,#2,#3\@nil{%  #1,#2 only for algebraic code
  \addbefore@par{algebraic=false}%
  \begin@SpecialObj
  \use@par
  \addto@pscode{
     gsave \psk@PSfont\space findfont \psk@fontscale scalefont setfont
     \ifPst@algebraic
       /x #1 def
       /Func (#2) tx@AlgToPs begin AlgToPs end cvx def
       Func
     \else #1 \fi
     \psk@decimals -1 gt { 10 \psk@decimals exp dup 3 1 roll mul cvi exch div } if
     %\psk@valuewidth 
     100 string cvs %/Output exch def % save output
     \ifPst@comma dot2comma \fi        % do we have to change dot to comma
     \psk@xShift\space 0 moveto  %Output
     show grestore
  }%
  \end@SpecialObj%
}

\define@boolkey[psset]{pst-func}[Pst@]{round}[true]{}%
\define@boolkey[psset]{pst-func}[Pst@]{science}[true]{%
  \ifPst@science\def\psk@Scin{true }\else\def\psk@Scin{false }\fi}
\psset[pst-func]{science=false,round=false}
\def\psPrintValueNew{\pst@object{psPrintValueNew}}
\def\psPrintValueNew@i#1{\expandafter\psPrintValueNew@ii#1,,\@nil}
\def\psPrintValueNew@ii#1,#2,#3\@nil{%  #1,#2 only for algebraic code
  \begin@SpecialObj
  \addto@pscode{  %		thanks to Buddy Ledger
     /mfont { \psk@PSfont\space findfont \psk@fontscale scalefont setfont } bind def
     /mfontexp { \psk@PSfont\space findfont \psk@fontscale 1.2 div scalefont setfont } bind def
     /s1 { /StandardSymbolsPS findfont \psk@fontscale scalefont setfont } bind def
     \ifPst@algebraic
        /x #1 def
        /Func (#2) tx@AlgToPs begin AlgToPs end cvx def
        Func
     \else #1 \fi
     /value ED
     \psk@Scin {
       value 0 ne { value log floor cvi /expon ED }{ /expon 0 def } ifelse
       value 10 expon exp div
       \psk@decimals -1 gt { 10 \psk@decimals exp dup 3 1 roll mul
         \ifPst@round round \else cvi \fi  exch div } if
       \psk@decimals 0 eq { cvi } if /numb ED
       expon \psk@valuewidth string cvs /expon exch def
       numb \psk@valuewidth string cvs
       \ifPst@comma dot2comma \fi  % do we have to change dot to comma
       /Output exch def
       /txspc \psk@fontscale 4 div def
       \psk@xShift\space 0 moveto mfont Output show
       txspc 0 rmoveto s1 (\string\264) show
       txspc 0 rmoveto mfont (10) show
       txspc 2 div txspc 1.5 mul rmoveto mfontexp expon show }
     { value
       \psk@decimals -1 gt { 10 \psk@decimals exp dup 3 1 roll mul
         \ifPst@round round \else cvi \fi exch div } if
       \psk@decimals 0 eq { cvi } if %inserted to handle decimals=0
       \psk@valuewidth string cvs
       \ifPst@comma dot2comma \fi         % do we have to change dot to comma
       \psk@xShift\space 0 moveto mfont %Output
       show
     } ifelse
  }%
  \end@SpecialObj%
}
%
% Integrals 2006-01-16
% Jose-Emilio Vila-Forcen, Herbert Voss
%
\def\psCumIntegral{\pst@object{psCumIntegral}}
\def\psCumIntegral@i#1#2#3{%
  \begin@OpenObj%
  \addto@pscode{
    /a #1  def
    /dx #2 #1 sub \psk@plotpoints\space div def
    /b a dx add def
    /scx { \pst@number\psxunit mul } def
    /scy { \pst@number\psyunit mul } def
    tx@FuncDict begin
     /SFunc \ifPst@algebraic (#3) tx@AlgToPs begin AlgToPs end cvx \else { #3 } \fi def
    end
    a scx 0 moveto
    \psk@plotpoints 1 sub {
      a b \psk@Simpson\space \ifPst@algebraic true \else false \fi  % a b M useX on Stack
      tx@FuncDict begin Simpson I end % y value on stack
      scy b scx exch lineto
      /b b dx add def
    } repeat
%    stroke
  }%
%  \psk@fillstyle%
%  \pst@stroke%
  \end@OpenObj%
}
%
\def\psIntegral{\pst@object{psIntegral}}
\def\psIntegral@i#1#2(#3,#4)#5{%
  \begin@OpenObj%
  \addto@pscode{
    /a #3  def
    /dx #4 #3 sub \psk@plotpoints\space div def
    /b #4 def
    /aa #1 def
    /dd #2 #1 sub \psk@plotpoints\space div def
    /t aa dd add def
    /scx { \pst@number\psxunit mul } def
    /scy { \pst@number\psyunit mul } def
    tx@FuncDict begin /SFunc { t #5 } def end
      a b \psk@Simpson false % a b M useX on Stack
      tx@FuncDict begin Simpson I end % y value on stack
      scy t scx exch moveto
      /t t dd add def
    \psk@plotpoints 1 sub {
      a b \psk@Simpson false % a b M useXon Stack
      tx@FuncDict begin Simpson I end % y value on stack
      scy t scx exch lineto
      /t t dd add def
    } repeat
%    stroke
  }%
%  \psk@fillstyle%
%  \pst@stroke%
  \end@OpenObj%
}
%
\def\psConv{\@ifnextchar[{\psConv@i}{\psConv@i[]}}
\def\psConv@i[#1]#2#3(#4,#5)#6#7{%
  \psIntegral[#1]{#2}{#3}(#4,#5){pop pop xVal #6\space xVal t neg add #7\space mul}%
}%
%
\define@key[psset]{pst-func}{PrintVLimit}[1e-6]{\def\psFunc@PrintVLimit{#1 }}
\define@key[psset]{pst-func}{Switch2Log}[80]{\def\psFunc@Switch2Log{#1 }}
\define@boolkey[psset]{pst-func}[Pst@]{printValue}[true]{}
\define@boolkey[psset]{pst-func}[Pst@]{LineEnding}[true]{}
\define@boolkey[psset]{pst-func}[Pst@]{VLines}[true]{}
\define@key[psset]{pst-func}{barwidth}[1]{\def\psFunc@barwidth{#1 }}% a factor, not a dimen
\define@key[psset]{pst-func}{rightEnd}[2]{\def\psFunc@rightEnd{#1 }}
\define@key[psset]{pst-func}{leftEnd}[2]{\def\psFunc@leftEnd{#1 }}
\define@key[psset]{pst-func}{labelangle}[90]{\def\psFunc@langle{#1 }}
\define@key[psset]{pst-func}{xlabelsep}[0]{\def\psFunc@xlabelsep{#1 }}
\define@key[psset]{pst-func}{radiusout}[2]{\def\psFunc@radiusout{#1 }}
\define@key[psset]{pst-func}{radiusinL}[0]{\def\psFunc@radiusinL{#1 }}
\define@key[psset]{pst-func}{radiusinR}[1.5]{\def\psFunc@radiusinR{#1 }}
\define@key[psset]{pst-func}{LabelColor}[black]{\pst@getcolor{#1}\psk@LabelColor}
\define@key[psset]{pst-func}{LineEndColorL}[green]{\pst@getcolor{#1}\psk@LineEndColorL}
\define@key[psset]{pst-func}{LineEndColorR}[red]{\pst@getcolor{#1}\psk@LineEndColorR}
\define@key[psset]{pst-func}{fillcolorA}[blue!40]{\pst@getcolor{#1}\psk@fillcolorA}
\define@key[psset]{pst-func}{fillcolorB}[red!40]{\pst@getcolor{#1}\psk@fillcolorB}
\define@key[psset]{pst-func}{vlinestyle}[solid]{\def\psvlinestyle{#1}}
\def\psfs@alternateColors{}
\def\pst@alternateColors{alternateColors}
%
\psset[pst-func]{printValue=false,barwidth=1,labelangle=90,xlabelsep=0,LabelColor=black,LineEndColorL=green,LineEndColorR=red,radiusout=2, radiusinL=0,radiusinR=1.5,LineEnding=true,leftEnd=1,rightEnd=2,VLines=false,
  fillcolorA=blue!40,fillcolorB=red!40,PrintVLimit=1e-64,Switch2Log=80}
%
\def\psBinomial{\pst@object{psBinomial}}
\def\psBinomial@i#1#2{\psBinomial@ii#1,,,\@nil{#2}}%
\def\psBinomial@ii#1,#2,#3,#4\@nil#5{%
  \def\pst@tempA{#2}%
  \ifx\pst@tempA\@empty
    \psBinomial@iii{0}{#1}{#1}{#5}%
  \else
    \def\pst@tempA{#3}%
    \ifx\pst@tempA\@empty\psBinomial@iii{#1}{#2}{#2}{#5}%
    \else\psBinomial@iii{#1}{#2}{#3}{#5}\fi
  \fi}
%
\def\psBinomial@iii#1#2#3#4{%
 \addbefore@par{valuewidth=15}%
  \begin@OpenObj
  \addto@pscode{%
    /toggle false def
    \ifx\psk@fillstylename\pst@alternateColors
       /ColA { \pst@usecolor\psk@fillcolorA } def
       /ColB { \pst@usecolor\psk@fillcolorB } def
    \fi
    /ValueSwitch2Log \psFunc@Switch2Log\space def
    /scx { \pst@number\psxunit mul } def
    /scy { \pst@number\psyunit mul } def
    /m #1 def
    /n #2 def
    /N #3 def
    /p #4 def
    /dx \psFunc@barwidth 2 div def
    /q 1 p sub def
    /kOld dx neg m add def
    kOld scx 0 moveto   % starting point#
%%----------------------------------------------------------------------------------------------------------------------------------------------------------------
%%-------------------- Code, der bei gro{\ss}em N \"{u}ber den Logarithmus arbeitet von Manuel Luque ---------------------------------------------
    /SumLog {%                                                                               log(1) + log(2) + ... + log(n)   =  log(n!)
      2 dict begin %                                                         Erwartet die Eingabe eines Wertes, n\"{a}mlich bis zu welchem Wert addiert werden soll ---> n
      /nV exch def %                                                       der \"{u}bergebene Wert wird in nV abgespeichert
      /iV 1 def %                                                               Startwert ist  iV=1
      0 nV { iV log add /iV iV 1 add def } repeat %  0 + log(1) + log(2) + ... + log(n)   Die Null ist dazu n\"{o}tig um die erste Addition machen zu k\"{o}nnen
      end
    } def
    /binomial_comb{ %                                                                      log(N!) - log(k!) - log((n-k)!)    =  log[N!/(k!*(n-k)!)]
      2 dict begin
      /nV exch def
      /NV exch def
      NV SumLog nV SumLog sub
      NV nV sub SumLog sub
      end
    } def
    /PnN {
      3 dict begin
      /pV exch def
      /nV exch def
      /NV exch def
      10
      NV nV binomial_comb
      nV pV log mul
      NV nV sub 1 pV sub log mul add add
      exp
      end
    } def
%%------------------------------------------------------------------------------------------------
%%------------------------------------------------------------------------------------------------
    N ValueSwitch2Log gt 
      { } 
      { 0 1 m 1 sub 
        {%
          /k exch def       % save loop variable
          k 0 eq
            { /Y q N exp def }
            { /Y Y N k sub 1 add mul k div p mul q div def }
          ifelse
        } for
      } ifelse
      m 1 n {%            % n-m+1 times
        /k exch def       % save loop variable
        N ValueSwitch2Log gt 
          { N k p PnN /Y exch def }
          { k 0 eq
            { /Y q N exp def }
            { /Y Y N k sub 1 add mul k div p mul q div def }
            ifelse 
          } 
          ifelse % recursive definition
      \ifx\psk@fillstylename\pst@alternateColors
        newpath
        kOld scx 0 moveto
      \fi
      kOld scx Y scy L k dx add scx Y scy L
      \ifPst@markZeros k dx add scx 0 L kOld 1 add scx 0 L \fi
      \ifx\psk@fillstylename\pst@alternateColors
        closepath
        gsave toggle {ColA}{ColB} ifelse \tx@setTransparency fill grestore
        /toggle toggle not def
      \fi
      \ifPst@printValue
        gsave \psk@PSfont\space findfont \psk@fontscale scalefont setfont \pst@usecolor\psk@LabelColor %
        Y \psFunc@PrintVLimit le { () } { Y \psk@valuewidth\space string cvs } ifelse
        \ifPst@comma dot2comma \fi
        k scx \psk@fontscale 2 div add \psFunc@xlabelsep scx add
        Y scy \pst@number\pslabelsep add moveto
        \psFunc@langle rotate show grestore
      \fi
      /kOld kOld 1 add def
    } for
    \ifPst@markZeros\else k dx add scx 0 L \fi % last line down to x-axis
  }%
%  \psk@fillstyle
%  \pst@stroke
  \end@OpenObj
}%
%
%
\def\psBinomialN{\pst@object{psBinomialN}}%
\def\psBinomialN@i#1#2{%
  \leavevmode
  \pst@killglue
  \begingroup
  \use@par
  \init@pscode
  \def\cplotstyle{curve}%
  \ifx\psplotstyle\cplotstyle \@nameuse{beginplot@\psplotstyle}\fi
  \addto@pscode{%
   \ifx\psplotstyle\cplotstyle /Curve true def \else /Curve false def \fi
    /scx { \pst@number\psxunit mul } def
    /scy { \pst@number\psyunit mul } def
    /N #1 def
    /p #2 def				% probability
    /q 1 p sub def
    /E N p mul def
    /sigma E q mul sqrt def 		% variant
    /dx 1.0 sigma div 2 div def
    /xOld dx neg E sub sigma div def
    /xEnd xOld neg dx add scx def
    Curve
      { /Coors [xOld dx sub scx 0] def }% saves the coordinates for curve
      { xOld scx 0 moveto }	% starting point
    ifelse
   0 1 N {%				% N times
      /k exch def			% save loop variable
      k 0 eq
        { /Y q N exp def }
        { /Y Y N k sub 1 add mul k div p mul q div def }
      ifelse % recursive definition
      /x k E sub sigma div dx add def
      /y Y sigma mul def		% normalize
      Curve
        { x dx sub scx y scy Coors aload length 2 add array astore /Coors exch def}
        { xOld scx y scy L x scx y scy L
          \ifPst@markZeros x scx 0 L \fi %
        } ifelse
      \ifPst@printValue
        gsave \psk@PSfont\space findfont \psk@fontscale scalefont setfont \pst@usecolor\psk@LabelColor %
        y \psk@valuewidth string cvs %/Output exch def
        \ifPst@comma dot2comma \fi      % do we have to change dot to comma
        x dx sub scx \psk@fontscale 2 div add \psFunc@xlabelsep scx add
        y scy \pst@number\pslabelsep add moveto
        \psFunc@langle rotate show grestore
      \fi
      /xOld x def
    } for
    Curve { [ xEnd 0 Coors aload pop } if% showpoints on top of the stack
  }%
  \ifx\psplotstyle\cplotstyle\@nameuse{endplot@\psplotstyle}\else
    \psk@fillstyle
    \pst@stroke
  \fi
  \use@pscode
  \endgroup
  \ignorespaces
}%
\def\psBinomialF{\pst@object{psBinomialF}}
\def\psBinomialF@i#1#2{\psBinomialF@ii#1,,,\@nil{#2}}%
\def\psBinomialF@ii#1,#2,#3,#4\@nil#5{%
  \def\pst@tempA{#2}%
  \ifx\pst@tempA\@empty
    \psBinomialF@iii{0}{#1}{#1}{#5}%
  \else
    \def\pst@tempA{#3}%
    \ifx\pst@tempA\@empty\psBinomialF@iii{#1}{#2}{#2}{#5}%
    \else\psBinomialF@iii{#1}{#2}{#3}{#5}\fi
  \fi}%
\def\psBinomialF@iii#1#2#3#4{%
  \addbefore@par{valuewidth=15}%
  \begin@OpenObj
  \addto@pscode{%
    /ValueSwitch2Log \psFunc@Switch2Log\space def
    /scx { \pst@number\psxunit mul } def
    /scy { \pst@number\psyunit mul } def
    /m #1 def
    /n #2 def
    /N #3 def
    /p #4 def
    /dx \psFunc@barwidth 2 div def
    /q 1 p sub def
    \ifPst@markZeros /kOld dx neg m add def \else
    /kOld m def \fi
    \psFunc@leftEnd 0 eq
      { }
      { kOld 0 eq
        { kOld \psFunc@leftEnd sub scx 0 moveto   % starting point
          \ifPst@LineEnding
            kOld \psFunc@radiusout\pst@number\psxunit div sub scx 0 L stroke
            kOld scx 0 newpath \psFunc@radiusout 0 360 arc kOld scx 0 
            \psFunc@radiusinR 360 0 arcn closepath \pst@usecolor\psk@LineEndColorR fill 
          \else
            kOld scx 0 L stroke 
          \fi 
        }{ } ifelse } ifelse
    kOld scx 0 moveto   % starting point
%%----------------------------------------------------------------------------------------------------------------------------------------------------------------
%%-------------------- Code, der bei gro{\ss}em N \"{u}ber den Logarithmus arbeitet von Manuel Luque ---------------------------------------------
    /SumLog {%                                                                               log(1) + log(2) + ... + log(n)   =  log(n!)
      2 dict begin %                                                         Erwartet die Eingabe eines Wertes, n\"{a}mlich bis zu welchem Wert addiert werden soll ---> n
      /nV exch def %                                                       der \"{u}bergebene Wert wird in nV abgespeichert
      /iV 1 def %                                                               Startwert ist  iV=1
      0 nV { iV log add /iV iV 1 add def } repeat %  0 + log(1) + log(2) + ... + log(n)   Die Null ist dazu n\"{o}tig um die erste Addition machen zu k\"{o}nnen
      end
    } def
    /binomial_comb{ %                                                                      log(N!) - log(k!) - log((n-k)!)    =  log[N!/(k!*(n-k)!)]
      2 dict begin
      /nV exch def
      /NV exch def
      NV SumLog nV SumLog sub
      NV nV sub SumLog sub
      end
    } def
    /PnN {
      3 dict begin
      /pV exch def
      /nV exch def
      /NV exch def
      10
      NV nV binomial_comb
      nV pV log mul
      NV nV sub 1 pV sub log mul add add
      exp
      end
    } def
%%------------------------------------------------------------------------------------------------
%%------------------------------------------------------------------------------------------------
    0 1 m 1 sub {
      /k exch def       % save loop variable
      N ValueSwitch2Log gt {  k 0 eq
        { N k p PnN /Y exch def /F Y def }
        { N k p PnN /Y exch def /F F Y add def }
        ifelse 
      }{ k 0 eq
         { /Y q N exp def /F Y def }
         { /Y Y N k sub 1 add mul k div p mul q div def
           /F F Y add def
         } ifelse
       } ifelse
     } for
    m 1 n {             % n-m+1 times
      /k exch def       % save loop variable
      N ValueSwitch2Log gt 
        { k 0 eq
          { N k p PnN /Y exch def /F Y def }
          { N k p PnN /Y exch def /F F Y add def } 
        ifelse }
        {  k 0 eq
          { /Y q N exp def /F Y def }
          { /Y Y N k sub 1 add mul k div p mul q div def /F F Y add def
        } ifelse
      } ifelse  % recursive definition
      F 1 gt { /F 1 def } if
      \ifPst@markZeros kOld scx F scy L k dx add scx F scy L k dx add scx 0 L kOld 1 add scx 0 L
      \else
       kOld scx F scy newpath \psFunc@radiusout 0 360 arc kOld scx F scy \psFunc@radiusinL 360 0 arcn closepath \pst@usecolor\psk@LineEndColorL fill
      kOld \ifPst@LineEnding\psFunc@radiusout\pst@number\psxunit div add \fi
      scx F scy moveto
      k 1 add
      \ifPst@LineEnding\psFunc@radiusout\pst@number\psxunit div sub \fi
       k N eq
      { \psFunc@rightEnd add }
      { } ifelse
      scx F scy L \pst@usecolor\pslinecolor [1 0] 0 setdash stroke
      \ifPst@LineEnding
      k N eq
      { }
      { k 1 add scx F scy newpath \psFunc@radiusout 0 360 arc k 1 add scx F scy \psFunc@radiusinR 360 0 arcn closepath \pst@usecolor\psk@LineEndColorR fill }
      ifelse
        \ifPst@VLines
        Y \psFunc@radiusout\pst@number\psyunit div 2 mul le
        { }
        { k scx F Y sub \psFunc@radiusout\pst@number\psyunit div add scy moveto k scx F \psFunc@radiusout\pst@number\psyunit div sub scy L \pst@usecolor\pslinecolor [3 2] 0 setdash stroke } ifelse \fi
      \else
      \ifPst@VLines
      Y \psFunc@radiusout\pst@number\psyunit div 0.1 mul le
      { }
      { k scx F Y sub scy moveto k scx F scy L \pst@usecolor\pslinecolor [3 2] 0 setdash stroke } ifelse
     \fi\fi\fi
      \ifPst@printValue
        gsave \psk@PSfont\space findfont \psk@fontscale scalefont setfont \pst@usecolor\psk@LabelColor %
        Y \psFunc@PrintVLimit le { () } { F \psk@valuewidth\space string cvs } ifelse
        \ifPst@comma dot2comma \fi
        k scx \psk@fontscale 2 div add \psFunc@xlabelsep scx add
        F scy \pst@number\pslabelsep add moveto
        \psFunc@langle rotate show grestore
      \fi
      /kOld kOld 1 add def
    } for
  }%
  \end@OpenObj%
}%
%
\def\psBinomialFS{\pst@object{psBinomialFS}}
\def\psBinomialFS@i#1#2{\psBinomialFS@ii#1,,,\@nil{#2}}%
\def\psBinomialFS@ii#1,#2,#3,#4\@nil#5{%
  \def\pst@tempA{#2}%
  \ifx\pst@tempA\@empty
    \psBinomialFS@iii{0}{#1}{#1}{#5}%
  \else
    \def\pst@tempA{#3}%
    \ifx\pst@tempA\@empty\psBinomialFS@iii{#1}{#2}{#2}{#5}%
    \else\psBinomialFS@iii{#1}{#2}{#3}{#5}\fi
  \fi}%
\def\psBinomialFS@iii#1#2#3#4{%
 \addbefore@par{valuewidth=15}%
  \begin@OpenObj
  \addto@pscode{%
    /ValueSwitch2Log \psFunc@Switch2Log\space def
    /scx { \pst@number\psxunit mul } def
    /scy { \pst@number\psyunit mul } def
    /m #1 def
    /n #2 def
    /N #3 def
    /p #4 def
    /dx \psFunc@barwidth 2 div def
    /q 1 p sub def
    \ifPst@markZeros /kOld dx neg m add def \else
    /kOld m def \fi
    kOld scx 0 moveto   % starting point
%%----------------------------------------------------------------------------------------------------------------------------------------------------------------
%%-------------------- Code, der bei gro{\ss}em N \"{u}ber den Logarithmus arbeitet von Manuel Luque ---------------------------------------------
    /SumLog {%                                                                               log(1) + log(2) + ... + log(n)   =  log(n!)
      2 dict begin %                                                         Erwartet die Eingabe eines Wertes, n\"{a}mlich bis zu welchem Wert addiert werden soll ---> n
      /nV exch def %                                                       der \"{u}bergebene Wert wird in nV abgespeichert
      /iV 1 def %                                                               Startwert ist  iV=1
      0 nV { iV log add /iV iV 1 add def } repeat %  0 + log(1) + log(2) + ... + log(n)   Die Null ist dazu n\"{o}tig um die erste Addition machen zu k\"{o}nnen
      end
    } def
    /binomial_comb{ %                                                                      log(N!) - log(k!) - log((n-k)!)    =  log[N!/(k!*(n-k)!)]
      2 dict begin
      /nV exch def
      /NV exch def
      NV SumLog nV SumLog sub
      NV nV sub SumLog sub
      end
    } def
    /PnN {
      3 dict begin
      /pV exch def
      /nV exch def
      /NV exch def
      10
      NV nV binomial_comb
      nV pV log mul
      NV nV sub 1 pV sub log mul add add
      exp
      end
    } def
%%------------------------------------------------------------------------------------------------
    0 1 m 1 sub {
     /k exch def       % save loop variable
     N ValueSwitch2Log gt {  k 0 eq
       { N k p PnN /Y exch def /F Y def /FS 1 F sub def }
       { N k p PnN /Y exch def /F F Y add def /FS 1 F sub def }
      ifelse }
         { k 0 eq
           { /Y q N exp def /F Y def /FS 1 F sub def }
           { /Y Y N k sub 1 add mul k div p mul q div def
             /F F Y add def /FS 1 F sub def }
         ifelse } ifelse
      } for
      m 1 n {             % n-m+1 times
        /k exch def       % save loop variable
        N ValueSwitch2Log gt { k 0 eq
          { N k p PnN /Y exch def /F Y def /FS 1 F sub def }
          { N k p PnN /Y exch def /F F Y add def /FS 1 F sub def } ifelse }
        {  k 0 eq
          { /Y q N exp def  /F Y def /FS 1 F sub def }
          { /Y Y N k sub 1 add mul k div p mul q div def
            /F F Y add def /FS 1 F sub def }
        ifelse
      } ifelse  % recursive definition
      FS 0 lt { /FS 0 def } if
      \ifPst@markZeros kOld scx FS scy L k dx add scx FS scy L k dx add scx 0 L kOld 1 add scx 0 L
      \else
      \ifPst@LineEnding
      k 0 eq
        { \psFunc@leftEnd neg scx 1 scy moveto \psFunc@radiusout neg 1 scy L gsave \pst@usecolor\pslinecolor stroke grestore } if
      k m eq {kOld scx \psFunc@radiusout add FS Y add scy moveto  kOld scx FS Y add scy  newpath \psFunc@radiusout 0 360 arc kOld scx FS Y add scy \psFunc@radiusinR 360 0 arcn closepath
      gsave \pst@usecolor\psk@LineEndColorR fill grestore } if
        kOld scx FS scy newpath \psFunc@radiusout 0 360 arc kOld scx FS scy \psFunc@radiusinL 360 0 arcn closepath \pst@usecolor\psk@LineEndColorL fill % } ifelse
       \fi
      kOld \ifPst@LineEnding\psFunc@radiusout\pst@number\psxunit div add \fi
%      k 0 eq
%      { \psFunc@leftEnd sub }
%      { } ifelse
      scx FS scy moveto
      k 1 add
      \ifPst@LineEnding\psFunc@radiusout\pst@number\psxunit div sub \fi
       k N eq
      { \psFunc@rightEnd add }
      { } ifelse
      scx FS scy L \pst@usecolor\pslinecolor [1 0] 0 setdash stroke
      \ifPst@LineEnding
      k N eq
      { }
      { k 1 add scx FS scy newpath \psFunc@radiusout 0 360 arc k 1 add scx FS scy \psFunc@radiusinR 360 0 arcn closepath \pst@usecolor\psk@LineEndColorR fill }
      ifelse
        \ifPst@VLines
        Y \psFunc@radiusout\pst@number\psyunit div 2 mul le
        { }
        {  k scx FS Y add \psFunc@radiusout\pst@number\psyunit div sub scy moveto k scx FS \psFunc@radiusout\pst@number\psyunit div add scy L \pst@usecolor\pslinecolor [3 2] 0 setdash stroke } ifelse \fi
      \else
      \ifPst@VLines
      Y \psFunc@radiusout\pst@number\psyunit div 0.1 mul le
      { }
      { k scx FS Y add scy moveto k scx FS scy L \pst@usecolor\pslinecolor [3 2] 0 setdash stroke } ifelse
     \fi\fi\fi
      \ifPst@printValue
        gsave \psk@PSfont\space findfont \psk@fontscale scalefont setfont \pst@usecolor\psk@LabelColor %
        Y \psFunc@PrintVLimit le { () } { FS \psk@valuewidth\space string cvs } ifelse
         \ifPst@comma dot2comma \fi
        k 0.1 add scx \psk@fontscale 2 div add \psFunc@xlabelsep scx add
        FS  scy \pst@number\pslabelsep add moveto
        \psFunc@langle rotate show grestore
      \fi
      /kOld kOld 1 add def
    } for
  }%
  \end@OpenObj
}%
%
%
\define@key[psset]{pst-func}{LeftClipX}[-1]{\def\psFunc@LeftClipX{#1 }}
\define@key[psset]{pst-func}{RightClipX}[-1]{\def\psFunc@RightClipX{#1 }}
\psset[pst-func]{LeftClipX=-1,RightClipX=1}
\def\psBinomialC{\pst@object{psBinomialC}}
\def\psBinomialC@i#1#2{%
  \leavevmode
  \pst@killglue
  \begingroup
    \use@par
  \begin{psclip}%
{\psframe[linestyle=none,fillstyle=none](\psFunc@LeftClipX,-0.1)(!#1 \psFunc@RightClipX\space add 1)}%
%  \use@par
  \init@pscode
  \beginplot@cspline
%    \addto@pscode{/D {} def mark}%
   \addto@pscode{%
    /scx { \pst@number\psxunit mul } def
    /scy { \pst@number\psyunit mul } def
    /N #1 def
    /p #2 def                                                         % Wahrscheinlichkeit
    /q 1 p sub def                                                % q = 1 - p
    /E N p mul def                                               % Erwartungswert
    /sigma E q mul sqrt def                              % Varianz
    /dx 1.0 2 div def
    /xOld dx neg 3 sub def
    /xEnd 0 def
    /Coors [xOld dx sub scx 0] def                % speichern der Koordinaten f\"{u}r den Spline
%%----------------------------------------------------------------------------------------------------------------------------------------------------------------
%%-------------------- Code, der \"{u}ber den Logarithmus arbeitet von Manuel Luque ---------------------------------------------
    /SumLog {%                        log(1) + log(2) + ... + log(n)   =  log(n!)
      2 dict begin %                  Erwartet die Eingabe eines Wertes, n\"{a}mlich bis zu welchem Wert addiert werden soll ---> n
      /nV exch def %                  der \"{u}bergebene Wert wird in nV abgespeichert
      /iV 1 def %                     Startwert ist  iV=1
      0 nV { iV log add /iV iV 1 add def } repeat %  0 + log(1) + log(2) + ... + log(n)   
      %Die Null ist dazu n\"{o}tig um die erste Addition machen zu k\"{o}nnen
      end
    } def
    /binomial_comb{ %                 log(N!) - log(k!) - log((n-k)!)    =  log[N!/(k!*(n-k)!)]
      2 dict begin
      /nV exch def
      /NV exch def
      NV SumLog nV SumLog sub
      NV nV sub SumLog sub
      end
    } def
    /PnN {
      3 dict begin
      /pV exch def
      /nV exch def
      /NV exch def
      10
      NV nV binomial_comb
      nV pV log mul
      NV nV sub 1 pV sub log mul add add
      exp
      end
    } def
%%------------------------------------------------------------------------------------------------
%%------------------------------------------------------------------------------------------------
%------------------------------- 3 Punkte um die Kurve links auf Null zu bringen mit dem weiter oben gew\"{a}hlten Startwert sind es also 4 ------------
    /x -2 dx add def
    /y 0 def		
    x dx sub scx y scy Coors aload length 2 add array astore /Coors exch def
    /x  -0.65 dx add def
    /y N 0 p PnN 10 -1.3 exp mul def		
    x dx sub scx y scy Coors aload length 2 add array astore /Coors exch def
    /x  -0.375 dx add def
    /y N 0 p PnN 3 div def		
    x dx sub scx y scy Coors aload length 2 add array astore /Coors exch def
%----------------------------  Ende der 3 Punkte und Beginn eigentlicher Code ----------------
    0 1 N {%            % Schleife geht von 0 bis N in 1-er Schritten
      /k exch def       % save loop variable
      N k p PnN /Y exch def
      /x k dx add def
      /y Y def		
      x dx sub scx y scy Coors aload length 2 add array astore /Coors exch def
      /xOld x def
    } for
%--------------------------- Ende des eigentlichen Codes und Beginn von  ---
%------------------------------- 4 Punkte um die Kurve rechts auf Null zu bringen -----------
    /x N 0.375 add dx add def
    /y Y 3 div def		
    x dx sub scx y scy Coors aload length 2 add array astore /Coors exch def
    /x N 0.65 add dx add def
    /y Y 10 -1.3 exp mul def		
    x dx sub scx y scy Coors aload length 2 add array astore /Coors exch def
    /x N 2 add dx add def
    /y 0 def		
    x dx sub scx y scy Coors aload length 2 add array astore /Coors exch def
    /x N 3 add dx add def
    /y 0 def		
    x dx sub scx y scy Coors aload length 2 add array astore /Coors exch def
%----------------------------  Ende der 4 Punkte und Laden der Koordinaten des arrays ; L\"{o}schen (pop) von showpoints??----------------------
    [ xEnd 1 add 0 Coors aload pop % showpoints on top of the stack
   }%
\endplot@cspline
\use@pscode
\end{psclip}%
\endgroup
\ignorespaces}%
%
\def\psPoisson{\pst@object{psPoisson}}% with contributions from Gerry Coombes
\def\psPoisson@i#1#2{\psPoisson@ii#1,,\@nil{#2}}%
\def\psPoisson@ii#1,#2,#3\@nil#4{%
  \def\pst@tempA{#2}%
  \ifx\pst@tempA\@empty\psPoisson@iii{0}{#1}{#4}\else
    \psPoisson@iii{#1}{#2}{#4}\fi}%
\def\psPoisson@iii#1#2#3{%  M N lambda
  \begin@OpenObj%
  \addto@pscode{
    /scx { \pst@number\psxunit mul } def
    /scy { \pst@number\psyunit mul } def
    /M #1 def
    /N #2 def
    /lambda #3 def
    /elambda Euler #3 neg exp def  % e^-lambda
    /dx \psFunc@barwidth 2 div def
    /kOld dx neg M add def         % addition of M here
    kOld scx 0 moveto              % starting point
    /Y elambda def                 % start value
    0 1 M 1 sub {                  % skip over first M-1 rectangles
      /k exch def                  % whilst recursing probabilities
      k 0 eq { /Y elambda def }{ /Y Y lambda mul k div def } ifelse
    } for                          % nothing happens if M=0
    M 1 N {                        % N-M+1 times
      /k exch def                  % save loop variable
      k 0 eq { /Y elambda def }{ /Y Y lambda mul k div def } ifelse
      kOld scx Y scy L k dx add scx Y scy L
      \ifPst@markZeros k dx add scx 0 L \fi
      \ifPst@printValue
        gsave \psk@PSfont\space findfont \psk@fontscale scalefont setfont
        Y \psk@valuewidth string cvs %/Output exch def
        \ifPst@comma dot2comma \fi         % do we have to change dot to comma
        k scx \psk@fontscale 2 div add
        Y scy \pst@number\pslabelsep add moveto
        90 rotate show grestore
      \fi
      /kOld kOld 1 add def
      \ifPst@markZeros kOld scx 0 moveto \fi
    } for
    \ifPst@markZeros \else k dx add scx 0 L \fi % last line down to x-axis
  }%
% \psk@fillstyle
% \pst@stroke
  \end@OpenObj%
}
%
\define@key[psset]{pst-func}{alpha}[0.5]{\pst@checknum{#1}\psk@alpha }
\define@key[psset]{pst-func}{beta}[0.5]{\pst@checknum{#1}\psk@beta }
\psset[pst-func]{alpha=0.5,beta=0.5}
%
\def\psGammaDist{\pst@object{psGammaDist}}
\def\psGammaDist@i#1#2{%
  \ifdim#1pt<\z@ \psframebox*{\color{red}!!!\#1 must be greater than 0!!!}
  \else
    \addbefore@par{plotpoints=500,alpha=0.5,beta=0.5}%
    \begin@OpenObj
    \psplot[algebraic=false]{#1}{#2}{
      \psk@beta x mul \psk@alpha exp x div Euler \psk@beta neg x mul \psk@alpha GAMMALN sub exp mul}
    \end@OpenObj%
  \fi%
  \ignorespaces%
}
%
\def\psBetaDist{\pst@object{psBetaDist}}
\def\psBetaDist@i#1#2{%
  \ifdim#1pt<\z@ \psframebox*{\color{red}!!!\#1 must be greater than 0!!!}
  \else
    \begingroup
    \addbefore@par{plotpoints=200,alpha=1,beta=1}%
    \addto@par{algebraic=false}%
%    \begin@OpenObj
    \use@par
    \psplot{#1}{#2}{
      \psk@beta \psk@alpha add GAMMA
      \psk@beta GAMMA \psk@alpha GAMMA mul div
      1 x sub \psk@beta 1.0 sub exp mul
      x \psk@alpha 1.0 sub exp mul }
%    \end@OpenObj%
    \endgroup
  \fi%
  \ignorespaces%
}
%
\def\psChiIIDist{\pst@object{psChiIIDist}}
\def\psChiIIDist@i#1#2{%
  \addbefore@par{plotpoints=500,nue=1}%
  \begin@OpenObj
%  \ifdim\psk@nue pt<\z@ \psframebox*{\color{red}!!!nue must be greater than 0!!!}
%  \else
     \psplot[algebraic=false]{#1}{#2}{%
      x 2 div \psk@nue 2 div exp x div Euler -0.5 x mul \psk@nue 2 div GAMMALN sub exp mul }%
%  \fi%
  \end@OpenObj%
  \ignorespaces%
}
%
\def\psTDist{\pst@object{psTDist}}
\def\psTDist@i#1#2{%
  \leavevmode
  \pst@killglue
  \begingroup
  \addbefore@par{plotpoints=500}%
  \use@par
  \ifdim\psk@nue pt<\z@ \psframebox*{\color{red}!!!nue must be greater than 0!!!}
  \else
     \psplot[algebraic=false]{#1}{#2}{
      1 x 2 exp \psk@nue div 1 add \psk@nue 1 add 2 div exp div
      \psk@nue Pi mul sqrt div
      Euler \psk@nue 1 add 2 div GAMMALN \psk@nue 2 div GAMMALN sub exp mul
     }%
  \fi%
  \endgroup%
  \ignorespaces%
}

\define@key[psset]{pst-func}{lambda}[1]{\def\psk@lambda{#1 }}
\psset[pst-func]{lambda=0.5}

\def\psNonCentralTDist{\pst@object{psNonCentralTDist}}
\def\psNonCentralTDist@i#1#2{%
  \leavevmode
  \pst@killglue
  \begingroup
  \addbefore@par{plotpoints=500}%
  \use@par
  \ifdim\psk@nue pt<\z@ \psframebox*{\color{red}!!!nue must be greater than 0!!!}
  \else
     \psplot[algebraic=false]{#1}{#2}{%   P(x) = a ( b/c + d/e )
       \psk@nue dup 2 div exp \psk@nue factorial mul  % nominator of a
       2 \psk@nue exp Euler \psk@lambda dup mul 2 div exp mul
           \psk@nue x dup mul add \psk@nue 2 div exp mul 
               \psk@nue 2 div GAMMA mul               % denominator of a
       div      % a
       2 sqrt \psk@lambda mul x mul 
          \psk@nue 2 div 1 add 
             1.5 
                \psk@lambda dup mul x dup mul mul 2 div \psk@nue x dup mul add div
          tx@FuncDict begin ConfHyperFunc end % Confluent Hypergeometric Funtion of the First Kind
          mul  % b
       \psk@nue x dup mul add \psk@nue 1 add 2 div GAMMA mul % c
       div % b/c
       \psk@nue 1 add 2 div 0.5 \psk@lambda dup mul x dup mul mul 2 div \psk@nue x dup mul add div
          tx@FuncDict begin ConfHyperFunc end  % d
       \psk@nue x dup mul add sqrt \psk@nue 2 div 1 add GAMMA mul % e
       div % d/e
       add % b/c + d/e
       mul % (b/c + d/e) * a
     }%
  \fi
  \endgroup
  \ignorespaces
}
%
\def\psFDist{\pst@object{psFDist}}
\def\psFDist@i#1#2{%
  \ifdim#1pt<\z@ \psframebox*{\color{red}!!!\#1 must be greater than 0!!!}
  \else
    \leavevmode
    \pst@killglue
    \begingroup
    \addbefore@par{plotpoints=500,mue=1}%
    \use@par
     \psplot[algebraic=false]{#1}{#2}{
      x \psk@mue mul \psk@nue div dup \psk@mue 2 div exp x div
      exch 1 add \psk@mue \psk@nue add 2 div exp div
      Euler \psk@mue \psk@nue add 2 div GAMMALN
      \psk@mue 2 div GAMMALN sub \psk@nue 2 div GAMMALN sub exp mul
     }%
    \endgroup%
  \fi%
  \ignorespaces%
}
%
\def\psNonCentralFDist{\pst@object{psNonCentralFDist}}
\def\psNonCentralFDist@i#1#2{%
  \ifdim#1pt<\z@ \psframebox*{\color{red}!!!\#1 must be greater than 0!!!}
  \else
    \leavevmode
    \pst@killglue
    \begingroup
    \addbefore@par{plotpoints=500,nue=1,mue=1,lambda=0,saveData=false}%
    \use@par
    \ifluatex
      \ifPst@saveData
        \luaexec{ require("pst-func")
           printDataPlotFileFDist(#1,#2,\psk@plotpoints,\psk@mue,\psk@nue,\psk@lambda,"\jobname-lua","true","true") }%
      \else
        \luaexec{ require("pst-func")
           printDataPlotFileFDist(#1,#2,\psk@plotpoints,\psk@mue,\psk@nue,\psk@lambda,"\jobname-lua","true","false") }%
      \fi
%      \typeout{\stringData}
      \pslistplot{\stringData}%
    \else
       \psplot[algebraic=false]{#1}{#2}[/F_BETA \psk@mue 2 div \psk@nue 2 div tx@FuncDict begin BETA end def]{
         /constA Euler \psk@lambda -2 div exp def
         /Sum 
           constA F_BETA
           div   %  first fraction
           \psk@mue \psk@nue div \psk@mue 2 div exp
           mul
           \psk@nue dup \psk@mue x mul add div \psk@mue \psk@nue add 2 div exp
           mul
           x \psk@mue 2 div 1 sub exp
           mul 
        def 
        /K 1 def
        {
          constA \psk@lambda 2 div K exp mul %  nominator 
          \psk@nue 2 div \psk@mue 2 div K add tx@FuncDict begin BETA end K factorial mul % denominator
          div   %  first fraction
          \psk@mue \psk@nue div \psk@mue 2 div K add exp
          mul
          \psk@nue dup \psk@mue x mul add div \psk@mue \psk@nue add 2 div K add exp
          mul
          x \psk@mue 2 div 1 sub K add exp
          mul 
          dup abs 1e-5 lt 
            { pop exit } 
            { Sum add /Sum ED 
              K 1 add /K ED } ifelse
            K 100 ge { (Error ConfHyperFunc: K=100) ==  exit } if
        } loop
        Sum 
      }%
    \fi
    \endgroup
  \fi
  \ignorespaces
}


%
\define@key[psset]{pst-func}{m}[0]{\def\psk@cauchy@m{#1 }}
\define@key[psset]{pst-func}{b}[1]{\def\psk@cauchy@b{#1 }}
\psset[pst-func]{m=0,b=1}
%
\def\psCauchy{\pst@object{psCauchy}}
\def\psCauchy@i#1#2{{%
    \pst@killglue%
    \addbefore@par{plotpoints=200}%
    \use@par%
    \psplot[algebraic=false]{#1}{#2}{
      \psk@cauchy@b dup dup mul x \psk@cauchy@m sub dup mul add div Pi div
    }%
}\ignorespaces}
%
\def\psCauchyI{\pst@object{psCauchyI}}
\def\psCauchyI@i#1#2{{%
    \pst@killglue%
    \addbefore@par{plotpoints=200}%
    \use@par%
    \psplot[algebraic=false]{#1}{#2}{
      x \psk@cauchy@m sub \psk@cauchy@b div ATAN1 DegtoRad Pi div 0.5 add
    }%
}\ignorespaces}
%
\def\psWeibull{\pst@object{psWeibull}}
\def\psWeibull@i#1#2{%
  \addbefore@par{plotpoints=500,alpha=1,beta=1}%
  \begin@OpenObj
  \def\pst@tempA{#1}%
  \ifdim#1pt<\z@ \psline(#1,0)(0,0)\def\pst@tempA{0}\fi
  \psplot[algebraic=false]{\pst@tempA}{#2}{
    \psk@alpha \psk@beta \psk@alpha neg Power mul % alpha*beta^(-alpha)
    x \psk@alpha 1 sub Power                      % x^(alpha-1)
    mul
    Euler x \psk@beta div \psk@alpha Power neg Power % e^(-(x/beta)^alpha))
    mul }
  \end@OpenObj%
  \ignorespaces%
}
\def\psWeibullI{\pst@object{psWeibullI}}
\def\psWeibullI@i#1#2{%
  \addbefore@par{plotpoints=500,alpha=1,beta=1}%
  \begin@OpenObj
  \def\pst@tempA{#1}%
  \ifdim#1pt<\z@ \psline(#1,0)(0,0)\def\pst@tempA{0}\fi
  \psplot[algebraic=false]{\pst@tempA}{#2}{
    1
    Euler x \psk@beta div \psk@alpha Power neg Power % e^(-(x/beta)^alpha))
    sub
  }%
  \end@OpenObj%
  \ignorespaces%
}
%
\define@key[psset]{pst-func}{pd}[0.22]{\pst@checknum{#1}\psk@probability }
\define@key[psset]{pst-func}{R2}[0.11]{\pst@checknum{#1}\psk@portfolio }
\psset[pst-func]{pd=0.22,R2=0.11}
%
\def\psVasicek{\pst@object{psVasicek}}
\def\psVasicek@i#1#2{%
  \addbefore@par{plotpoints=500}%
  \begin@OpenObj
  \psplot{#1}{#2}[/pd \psk@probability\space def /R2 \psk@portfolio\space def ]{x tx@FuncDict begin vasicek end}
  \end@OpenObj%
  \ignorespaces%
}
\define@boolkey[psset]{pst-func}[Pst@]{Gini}[true]{}
\psset[pst-func]{Gini=false}
%
\def\psLorenz{\def\pst@par{}\pst@object{psLorenz}}
\def\psLorenz@i{\@ifnextchar<\psLorenz@ii{\psLorenz@ii<>}}
\def\psLorenz@ii<#1>#2{%
  \if@star\addto@par{fillstyle=solid,fillcolor=\pslinecolor}\fi
  \begingroup
  \use@par
  \def\Lorenz@code{
    25 dict begin
    [ #2 ] dup length /nVal ED 
    bubblesort /Yval ED
    Yval arraySum /Ysum ED
    Yval { Ysum div } forall nVal array astore /YvalRel ED %  Pred dist
    \if$#1$
      /temp {1 nVal div} bind def
      [ nVal { temp } repeat ] /XvalRel ED
      [ 1 1 nVal { nVal div } for ] /XvalRelCum ED
      %YvalRel bubblesort /YvalRel ED      
    \else
      [ #1 ] /Xval ED
      Xval arraySum /Xsum ED
      Xval { Xsum div } forall nVal array astore /XvalRel ED %  Freq dist
      /CumFreq 0 def
      [
        0 1 nVal 1 sub {
          XvalRel exch get CumFreq add dup /CumFreq ED
        } for
      ] /XvalRelCum ED
%      XvalRelCum ==
    \fi
    /CumPred 0 def
    [
    0 1 nVal 1 sub {
      YvalRel exch get CumPred add dup /CumPred ED
    } for
    ] /YvalRelCum ED
%    YvalRelCum ==
%    Xsum == Ysum ==
    \ifPst@Gini
      [ 0 1 nVal 1 sub { 
        /Index ED
        YvalRelCum Index get 
        Index 0 eq { 0 }{ YvalRelCum Index 1 sub get } ifelse
        add 0.5 mul XvalRel Index get mul 
      } for
      ] /Area ED
%      Area ==
      Area arraySum /B ED
      0.5 B sub dup B add div
      \psk@PSfont\space findfont \psk@fontscale scalefont setfont
      \psk@decimals -1 gt { 10 \psk@decimals exp dup 3 1 roll mul cvi exch div } if
      \psk@valuewidth string cvs %/Output exch def % save output
      \ifPst@comma dot2comma \fi      % do we have to change dot to comma
      /Output ED
      \psk@xShift\space -30 moveto (Gini: ) show
      Output show
    \fi
    0 0 % the start point
    0 1 nVal 1 sub { dup XvalRelCum exch get exch YvalRelCum exch get } for
%    pstack
    \if@star 1 0 0 0 \fi                   % add values for the closed curve
    end
  }% filling the area under the curve.
  \if@star\listplot*{\Lorenz@code}\else\listplot{\Lorenz@code}%
%     \listplot[plotstyle=bezier,linecolor=red]{\Lorenz@code}
  \fi
  \endgroup
  \ignorespaces}
%
% Superellipese / Lamefunction
\define@key[psset]{pst-func}{radiusA}[1]{\pst@getlength{#1}\pst@radiusA}
\define@key[psset]{pst-func}{radiusB}[1]{\pst@getlength{#1}\pst@radiusB}
\define@key[psset]{pst-func}{endAngle}[0]{\pst@getangle{#1}\psk@endAngle}
\psset[pst-func]{radiusA=1,radiusB=1,endAngle=0}%  startAngle is defined in PSTricks
%
\def\psLame{\pst@object{psLame}}
\def\psLame@i#1{%
  \leavevmode
  \pst@killglue
  \begingroup
  \addbefore@par{plotpoints=200,startAngle=0,endAngle=360,algebraic=false}%
  \use@par
  \parametricplot{\psk@startAngle}{\psk@endAngle}{%
     t cos dup mul 1 #1\space div exp \pst@radiusA \pst@number\psxunit div mul
     t 90 gt { t 270 lt { neg } if } if
     t sin dup mul 1 #1\space div exp \pst@radiusB \pst@number\psyunit div mul
     t 180 gt { neg } if }
  \endgroup
  \ignorespaces
}
%
\def\psWeierstrass{\pst@object{psWeierstrass}}
\def\psWeierstrass@i(#1,#2){\@ifnextchar[{\psWeierstrass@ii(#1,#2)}{\psWeierstrass@iii(#1,#2)}}
\def\psWeierstrass@ii(#1,#2)[#3]#4{% #1 xStart; #2 xEnd; #3 a; #4 b  %% the original definition
  \addbefore@par{plotpoints=500}%
  \begin@SpecialObj
  \psplot{#1}{#2}[/ps@a #3 def /ps@b #4 def]{
    /ps@n 0 def
    /ps@WS 0 def
    {
      ps@a ps@n exp ps@b ps@n exp Pi mul x mul RadtoDeg cos mul
      dup abs \psk@epsilon gt { ps@WS add /ps@WS ED }{ ps@WS add exit } ifelse
      ps@n 1 add /ps@n ED
    } loop
  }%
  \end@SpecialObj
}
\def\psWeierstrass@iii(#1,#2)#3{% #1 xStart; #2 xEnd; #3 a
  \addbefore@par{plotpoints=500}%
  \begin@SpecialObj
  \psplot{#1}{#2}[/ps@a #3 def]{
    /ps@k 1 def
    /ps@WS 0 def
    {
      ps@k ps@a exp Pi mul dup x mul RadtoDeg sin exch div
      dup abs \psk@epsilon gt { ps@WS add /ps@WS ED }{ ps@WS add exit } ifelse
      ps@k 1 add /ps@k ED
    } loop
  }%
  \end@SpecialObj
}
%
% For polar plots
%\define@boolkey[psset]{pst-func}[PstAdd@]{polarplot}[true]{}
%\psset[pst-func]{polarplot=false}
%
%\define@boolkey[psset]{pstricks-add}[Pst@]{GetFinalState}[true]{}
%\define@key[psset]{pstricks-add}{filename}{\def\psk@filename{#1}}%
%\define@boolkey[psset]{pstricks-add}[Pst@]{saveData}[true]{} % \ifPst@saveData
%\psset[pstricks-add]{GetFinalState=false,saveData=false,filename=PSTdata}
%
\define@key[psset]{pst-func}{stepFactor}[0.67]{\pst@checknum{#1}\psk@stepFactor }
\psset[pst-func]{stepFactor=0.67}
%
\def\psplotImp{\pst@object{psplotImp}}% 20060420
\def\psplotImp@i(#1,#2)(#3,#4){%
  \@ifnextchar[{\psplotImp@ii(#1,#2)(#3,#4)}{\psplotImp@ii(#1,#2)(#3,#4)[]}}
\def\psplotImp@ii(#1,#2)(#3,#4)[#5]#6{%
  \addbefore@par{filename=\jobname.data}%
  \begin@OpenObj%
  \addto@pscode{
    \ifPst@saveData /Pst@data (\psk@filename) (w) file def \fi
    /xMin #1 def
    /xMax #3 def
    /yMin #2 def
    /yMax #4 def
    #5 % additional PS code
    \ifPst@polarplot
      /@PolarAlgPlot (#6) tx@addDict begin AlgParser end cvx def
      /Func {
        /phi y x atan def
        /r x y Pyth def
        \ifPst@algebraic @PolarAlgPlot \else #6 \fi } def
    \else
      /Func \ifPst@algebraic (#6) tx@addDict begin AlgParser end cvx \else { #6 } \fi  def
    \fi
    /xPixel xMax xMin sub \pst@number\psxunit mul round cvi def
    /yPixel yMax yMin sub \pst@number\psyunit mul round cvi def
    /dx xMax xMin sub xPixel div def
    /dy yMax yMin sub yPixel div def
    /setpixel {
      dy div exch
      dx div exch
      \ifPst@saveData
        2 copy
        \pst@number\psyunit div exch  \pst@number\psxunit div
        20 string cvs Pst@data exch writestring
        Pst@data (\space) writestring
        20 string cvs Pst@data exch writestring
%        Pst@data (\string\]) writestring
        Pst@data (\string\n) writestring
      \fi
      \pst@number\pslinewidth 2 div 0 360 arc fill } bind def
%
    /VZ true def % suppose that F(x,y)>=0
    /x xMin def /y yMin def Func 0.0 lt { /VZ false def } if % erster Wert
    xMin dx \psk@stepFactor\space mul xMax {
      /x exch def
      \ifPst@saveData Pst@data ([\string\n) writestring \fi
      yMin dy \psk@stepFactor\space mul yMax {
	/y exch def
	Func 0 lt
	    { VZ { x y setpixel /VZ false def} if }
	    { VZ {}{ x y setpixel /VZ true def } ifelse } ifelse
      } for
      \ifPst@saveData Pst@data (]\string\n) writestring \fi
    } for
%% the same for the other way round without saving the data
    /VZ true def % suppose that F(x,y)>=0
    /x xMin def /y yMin def Func 0.0 lt { /VZ false def } if % erster Wert
    yMin dy \psk@stepFactor\space mul yMax {
      /y exch def
      \ifPst@saveData Pst@data ([\string\n) writestring \fi
      xMin dx \psk@stepFactor\space mul xMax {
	/x exch def
	Func 0 lt
	    { VZ { x y setpixel /VZ false def} if }
	    { VZ {}{ x y setpixel /VZ true def } ifelse } ifelse
      } for
      \ifPst@saveData Pst@data (]\string\n) writestring \fi
    } for
%
\iffalse
    /x xMin def /y yMin def Func 0.0 lt { /VZ false def } if % erster Wert
    yMin dy \psk@stepFactor\space mul yMax {
      /y exch def
      xMin dx \psk@stepFactor\space mul xMax {
	/x exch def
	Func 0 lt
	    { VZ { x y setpixel /VZ false def} if }
	    { VZ {}{ x y setpixel /VZ true def } ifelse } ifelse
      } for
    } for
\fi
    \ifPst@saveData Pst@data closefile \fi
  }%
  \end@OpenObj%
}
%
\def\psVolume{\pst@object{psVolume}}% 2007-06-23
\def\psVolume@i(#1,#2)#3#4{%
  \leavevmode
  \pst@killglue
  \begingroup
  \use@par
  \psplot[fillstyle=none]{#1}{#2}{#4}%     original function
  \ifPst@algebraic
    \psplot[fillstyle=none]{#1}{#2}{-(#4)}
  \else
    \psplot[fillstyle=none]{#1}{#2}{#4 neg}% mirrored at the x-axis
  \fi
  \multido{\iA=1+1}{#3}{%	 run it #3 times with increment \A
    \pscustom{%                  to get a closed filled ellipse
      \code{ %					 the PS code
        /dX #2 #1 sub #3 div def % delta x, the step
        /Start dX \iA\space 1 sub mul #1 add def % xStart
        /End Start dX add def                    % xEnd=xStart+dX
        /Height End Start add 2 div /x ED
          \ifPst@algebraic (#4) AlgParser cvx exec \else #4 \fi def } % height=f(x)
         % x is the mean between Start+End
      \psellipticarc(!Start 0)(! Height 8 div Height){90}{270}
      % draw the first falf of the ellipse
      \rlineto(! dX 0)% draw a line in x-direction
      \psellipticarc(!End 0)(! Height 8 div Height){270}{90}
      % draw the other half of the ellipse
      \rlineto(!dX neg 0)}}% draw a line in negative x-direction
  \psset{fillstyle=none}
%  \psellipse(#2,0)(!#2 dup #1 sub #3 div 2 div sub /x ED #4 dup
%	 8 div exch)% draw again the ellipse to get the borderline.
  \psellipse(#2,0)(!{ #2 dup #1 sub #3 div 2 div sub /x ED
         \ifPst@algebraic (#4) AlgParser cvx exec \else #4 \fi dup
	 8 div exch})  % draw again the ellipse to get the borderline.
  \psset{plotstyle=line,linestyle=dashed,
	 plotpoints=40,dotstyle=*,dotsize=0.5pt}
  \psplot[fillstyle=none]{#1}{#2}{#4}
  \ifPst@algebraic
    \psplot[fillstyle=none]{#1}{#2}{-(#4)}
  \else
    \psplot[fillstyle=none]{#1}{#2}{#4 neg}% mirrored at the x-axis
  \fi
  % draw again the curves to get the borderline
  \endgroup%
  \ignorespaces%
}
\def\txFunc@BezierCurve{ tx@FuncDict begin BezierCurve Points end }
\def\txFunc@BezierShowPoints{ tx@Dict begin /Points ED BezierShowPoints end }
\def\pst@BezierType{2 }	% the default
%
\def\psBezier#1{%		% allowed order is 1 ... 9
  \ifnum#1>0 \ifnum#1<10 \def\pst@BezierType{#1 }\fi\fi%
  \pst@object{psBezier}}
\def\psBezier@i{%
  \pst@getarrows{%
    \addbefore@par{plotpoints=200}%
    \begin@OpenObj
    \pst@getcoors[\psBezier@ii%
}}
\def\psBezier@ii{%
  \addto@pscode{%
    \psk@plotpoints	    	% step for Bezier T=0,0+epsilon,0+i*epsilon,...,1
    \pst@BezierType 		% type of the Bezier curve 2,3,4,...
    \txFunc@BezierCurve
    \ifshowpoints \txFunc@BezierShowPoints \else pop \fi
  }%
  \end@OpenObj}
%
\def\tx@Bernstein{ tx@FuncDict begin Bernstein end }
\define@boolkey[psset]{pst-func}[Pst@]{envelope}[true]{}
\psset[pst-func]{envelope=false}
%
\def\psBernstein{\pst@object{psBernstein}}%  \psBernstein[options](t1,t2)(i,n)
\def\psBernstein@i(#1,#2){%
  \@ifnextchar({\psBernstein@ii(#1,#2)}{\psBernstein@ii(0,1)(#1,#2)}}
%
\def\psBernstein@ii(#1,#2)(#3,#4){%	(tStart,tEnd)(i,n)
  \addbefore@par{plotpoints=200}%
  \begin@OpenObj
  \addto@pscode{%
    /ScreenCoor { \tx@ScreenCoor } def
    #1\space #2\space
    1.0 \psk@plotpoints\space div	    	% step=1/plotpoints
    #3\space #4\space
%   on stack we have tStart tEnd epsilon i n
    \ifPst@envelope true \else false \fi
    \tx@Bernstein
  }%
  \end@OpenObj}
%
\def\psThomae{\pst@object{psThomae}}
\def\psThomae@i(#1,#2)#3{%
  \addbefore@par{dotsize=1pt}
  \begin@ClosedObj
  \addto@pscode{
    1 1 #3 {
      dup
      /ipSave ED	% save loop value
      /ip ED		% dito
      1 1 #3 {		
        dup		
        /iqSave ED	% save loop value
        /iq ED		% dito
        {
          iq 0 le { exit } if
          ip iq mod
          /ip iq def
          /iq ED
        } loop
        ip 1 eq {
          /xVal ipSave iqSave div def
          xVal #1 ge { xVal #2 le {
            \psk@dotsize
            \@nameuse{psds@\psk@dotstyle}
            \pst@usecolor\pslinecolor xVal 1 iqSave div \tx@ScreenCoor
            2 copy pop 0 moveto lineto } if } if
%            2 copy moveto Dot } if } if
        } if
      } for
    } for
  }%
  \end@ClosedObj%
}
%
\def\psCplot{\def\pst@par{}\pst@object{psCplot}}
\def\psCplot@i#1#2#3#4{% start | end | complex variables | function
  \pst@killglue
  \begingroup
    \use@par
    \@nameuse{beginplot@\psplotstyle}%
    \addto@pscode{%
      \psplot@init
      /x #1 def
      /x1 #2 def
      /dx x1 x sub \psk@plotpoints div def
      #3
      /xy {
%        x
        tx@FuncDict begin
        #4 aload pop \pst@number\psyunit mul exch \pst@number\psxunit mul exch
        end
      } def}%
    \gdef\psplot@init{}%
    \@pstfalse
    \@nameuse{testqp@\psplotstyle}%
    \if@pst
      \psplot@ii
    \else
      \psplot@iii
    \fi
  \endgroup
  \ignorespaces}
%
\define@boolkey[psset]{pst-func}[Pst@]{xory}[true]{}
\define@boolkey[psset]{pst-func}[Pst@]{approx}[true]{}
\define@boolkey[psset]{pst-func}[Pst@]{Framed}[true]{}
\define@boolkey[psset]{pst-func}[Pst@]{Newton}[true]{}
\define@boolkey[psset]{pst-func}[Pst@]{PrintCoord}[true]{}
\define@boolkey[psset]{pst-func}[Pst@]{onlyNode}[true]{}
\define@boolkey[psset]{pst-func}[Pst@]{onlyYVal}[true]{}
\define@boolkey[psset]{pst-func}[Pst@]{originV}[true]{}
\define@key[psset]{pst-func}{PointName}[I]{\def\psk@func@PointName{#1}}
\define@key[psset]{pst-func}{ydecimals}[2]{\pst@getint{#1}\psk@ydecimals }
\define@key[psset]{pst-func}{labeldistance}[0]{\def\psk@ldistance{#1 }}
\psset[pst-func]{originV=false,onlyNode=false,ydecimals=2,labeldistance=0,
  PrintCoord=false,onlyYVal=false,Newton=false,PointName=I,Framed=false,xory=false,approx=true}
%
%--------------------------------------------------------------------
%------------- calculate the value of an intersectionpoint -----------
%---------------------------------------------------------------------
%
\def\psZero{\def\pst@par{}\pst@object{psZero}}
\def\psZero@i(#1,#2)#3{\@ifnextchar[{\psZero@ii(#1,#2)#3}{\psZero@ii(#1,#2){#3}[0]}}
\def\psZero@ii(#1,#2)#3[#4]#5{%
% (#1,#2) Intervall f\"{u}r die Nullstelle, bzw #1 Startwert f\"{u}r Newton, #3 1. Funktion, #4 2. Funktion, #5 Knotenname
 \begingroup
 \pst@killglue
 \addbefore@par{labelangle=0,framesep=1pt}%fontscale=40,PSfont=NimbusRomNo9L-Regu}%
 % labelangle=0 -> take xShift and yShift into account, otherwise xShift=pslabelsep * cos(labelangle)
 %                                                                yShift=pslabelsep * sin(labelangle)
 \use@par
 \pst@Verb{
   /FunctionA
     \ifPst@algebraic (#3) tx@AlgToPs begin AlgToPs end cvx \else {#3} \fi def
   /FunctionB
     \ifPst@algebraic (#4) tx@AlgToPs begin AlgToPs end cvx \else {#4} \fi def
 \ifPst@Newton %   ---------------------------        Newton-Verfahren ---  TS 05.2014
  /eps 1e-6 def  %                                         eps f\"{u}r (f(x+eps)-f(x))/eps
  /x0 #1 def     %                                         Startwert wird vorgegeben
  /Xend x0 def
   {
  /xe x0 eps add def %                                      xe=x+eps
  /x x0 def          %                                      x nun x0 setzen, damit f(x0) berechnet werden kann
  /func0 FunctionA FunctionB sub def % func0=F(x0)-FB(x0)
  func0 0 eq {exit} if %                                    wenn func0 = 0 -> exit
  /x xe def  %                                              x xe setzen
  /func1 FunctionA FunctionB sub def % func1=F(xe)-FB(xe)
  /Xend x0 func0 eps mul func1 func0 sub div sub def % Xend=x0-(func0*eps)/(func1-func0)
  /x Xend def %                                             x wird nun Xend
  /funcend FunctionA FunctionB sub def % funcend=F(Xend)-FB(Xend)
  funcend abs 1e-6 le {exit} if %                           Wenn funcend < 1e-6 -> exit
  /x0 Xend def %                                            sonst x0 = Xend und
  } loop %                                                  die Schleife beginnt mit x0 = Xend
  /xM Xend def
 \else  % --------------------   Intervallhalbierungsverfahren von Manuel L.
  /Xinf #1 def %                                            linker Wert des Intervalls, in dem die Nullstelle gesucht wird
  /Xsup #2 def %                                            linker Wert des Intervalls, in dem die Nullstelle gesucht wird
  { /xM Xinf Xsup add 2 div def %                           Mittelwert -> Mitte des Intervalls ist xM2
  /x Xinf def  %                                            Variable x wird als linker Wert festgelegt
  /F_1 FunctionA FunctionB sub def %                         Definition der Differenzfunktion F_1 an der Stelle x=Xinf2
  /x xM def  %                                              Variable x wird als Mittelwert festgelegt
  /F_M FunctionA FunctionB sub def %                         Definition der Differenzfunktion F_M an der Stelle x=xM2
  F_M 0 eq {exit} if %                                      Ist F_M = 0 => fertig
  F_1 F_M mul 0 ge {/Xinf xM def} {/Xsup xM def} ifelse %   F_1 * F_M > 0 (F_1 und F_M haben gleiches VZ) => neuer linker x-Wert = xM2
  Xinf Xsup sub abs 1e-6 le {exit} if } loop %              Die Schleife endet, wenn die Differenz von linkem und rechtem Wert < 10^-6
 \fi
  /x xM def FunctionA /yM exch def %                               x wird Mittelwert und an dieser Stelle die Funktion berechnet und als yM definiert
 }%
\pnode(! xM yM){#5}%
\addto@pscode{%
  \iftrue %\psFunc@langle pt>\z@
    /xShift \psFunc@langle cos \psk@ldistance mul def 
    /yShift \psFunc@langle sin \psk@ldistance mul def 
  \else
    /xShift \psk@xShift\space def
    /yShift \psk@yShift\space def
  \fi
\ifPst@onlyNode\else
  /dec \psk@decimals\space def   % Anzahl der Dezimahlen f\"{u}r x
  /ydec \psk@ydecimals\space def % Anzahl der Dezimalen f\"{u}r y
  /symb { /StandardSymbolsPS findfont \psk@fontscale\space scalefont setfont } bind def             % Symbolschrift f\"{u}r Klammer und senkrechten Strich
  /schrift { \psk@PSfont\space findfont \psk@fontscale scalefont setfont } bind def            % Schrift f\"{u}r die Ergebnisse
  /TiefSchrift { /StandardSymbolsPS findfont \psk@fontscale 1.8 div scalefont setfont } bind def    % Schrift f\"{u}r tiefergestellte Buchstaben
  /spdx \psk@fontscale 4 div def % Abstand, der mit der Schriftg\"{o}{\ss}e skaliert ist
  /spdy \psk@fontscale 15 div def % Abstand, der mit der Schriftg\"{o}{\ss}e skaliert ist
  /Wert { dec -1 le { /dec 15 def } if 10 dec exp mul round 10 dec exp div dec 0 eq { cvi 15 string cvs } { 15 string cvs } ifelse % x-Wert runden und als string ablegen
  \ifPst@comma dot2comma \fi show } def % gegebenenfalls Komma statt Punkt und Ergebnis anzeigen
  /yWert { 10 ydec exp mul round 10 ydec exp div ydec 0 eq { cvi 15 string cvs } { 15 string cvs } ifelse % y-Wert runden und als string ablegen
     \ifPst@comma dot2comma \fi show } def % gegebenenfalls Komma statt Punkt und Ergebnis anzeigen
  /FunctionA \ifPst@algebraic (#3) tx@AlgToPs begin AlgToPs end cvx \else {#3} \fi def
  /FunctionB \ifPst@algebraic (#4) tx@AlgToPs begin AlgToPs end cvx \else {#4} \fi def
  /frs \pst@number\psframesep def
 \ifPst@Newton
  /eps 1e-6 def
  /x0 #1 def
  /Xend x0 def
  { /xe x0 eps add def
    /x x0 def
    /func0 FunctionA FunctionB sub def
    func0 0 eq {exit} if
    /x xe def
    /func1 FunctionA FunctionB sub def
    /Xend x0 func0 eps mul func1 func0 sub div sub def
    /x Xend def
    /funcend FunctionA FunctionB sub def
    funcend abs 1e-6 le {exit} if
    /x0 Xend def
  } loop
  /xM Xend def
 \else
  /Xinf #1 def   /Xsup #2 def
  { /xM Xinf Xsup add 2 div def
  /x Xinf def  /F_1 FunctionA FunctionB sub def
  /x Xsup def  /F_2 FunctionA FunctionB sub def
  /x xM def    /F_M FunctionA FunctionB sub def
  F_M 0 eq {exit} if
  F_1 F_M mul 0 ge {/Xinf xM def} {/Xsup xM def} ifelse
  Xinf Xsup sub abs 1e-6 le {exit} if } loop
 \fi
  /x xM def FunctionA /yM exch def
  \ifPst@originV 0 0 \else
    xM xShift add \pst@number\psxunit mul
    yM yShift add \pst@number\psyunit mul
  \fi
  \pst@usecolor\pslinecolor
 \ifPst@PrintCoord
/xytext { moveto schrift (\psk@func@PointName) show
    0 spdy 2 mul neg rmoveto TiefSchrift (\psk@postString) show
    0 spdy 3 mul rmoveto symb (\string\050) show
    0 spdy neg rmoveto schrift xM Wert
    spdx 3 div spdy rmoveto symb (\string\174) show
    spdx 3 div spdy neg rmoveto schrift yM yWert 0 spdy rmoveto symb (\string\051) show } bind def
    xytext
    0 spdy neg rmoveto
% ------------------------------------------------------------------------------ TS 03.2017
    \ifPst@Framed
       frs \psk@fontscale 0.15 mul frs add neg rmoveto
       0 \psk@fontscale 0.9 mul frs 2 mul add rlineto
       \ifPst@originV
          frs neg \psk@fontscale 0.75 mul frs add L
         \else
          xM xShift add \pst@number\psxunit mul frs sub
          yM yShift add \pst@number\psyunit mul \psk@fontscale 0.75 mul add frs add L
       \fi
        0 \psk@fontscale 0.9 mul frs 2 mul add neg rlineto closepath
         gsave
          \pst@usecolor\psfillcolor \tx@setTransparency fill
         grestore
            \ifx\pslinestyle\@none\else
            \pst@number\pslinewidth SLW \pst@usecolor\pslinecolor \tx@setStrokeTransparency \@nameuse{psls@\pslinestyle} stroke
           \fi
       \ifPst@originV 0 0 \else
          xM xShift add \pst@number\psxunit mul
          yM yShift add \pst@number\psyunit mul
       \fi
         xytext
         0 spdy neg rmoveto
    \fi
%------------------------------ Wenn nicht die Form P(a|b) sondern x = a oder y = b oder nur a oder nur b -------------------
  \else
    \ifPst@originV 0 0 \else
       xM xShift add \pst@number\psxunit mul
       yM yShift add \pst@number\psyunit mul
    \fi
      /xorytext { moveto schrift
  \ifPst@onlyYVal%  nur der y-Wert
        \ifPst@xory%   in Form von y=b
          (\psk@func@PointName) show
          spdx 5 div neg spdy 2.2 mul neg rmoveto TiefSchrift (\psk@postString) show
          spdx 1.5 div spdy 2 mul rmoveto symb(\ifPst@approx \string \273 \else \string = \fi) show
          spdx 1.5 div neg 0 rmoveto schrift yM yWert
        \else schrift yM yWert
       \fi
     \else
       \ifPst@xory%     in Form x = a
          (\psk@func@PointName) show
          spdx 7 div neg spdy 2.2 mul neg rmoveto TiefSchrift (\psk@postString) show
          spdx 1.5 div spdy 2 mul rmoveto symb (\ifPst@approx \string \273 \else \string = \fi) show
          spdx 1.5 div neg 0 rmoveto schrift xM Wert
        \else schrift xM Wert
       \fi
     \fi } bind def
   xorytext
   \ifPst@Framed%------------------------------  Rahmenbeginn ------------------------------
       frs \psk@fontscale 0.14 mul frs add neg rmoveto
       0 \psk@fontscale 0.9 mul frs 2 mul add rlineto
          \ifPst@originV
             frs neg \psk@fontscale 0.75 mul frs add L
           \else
             xM xShift add \pst@number\psxunit mul frs sub
             yM yShift add \pst@number\psyunit mul \psk@fontscale 0.75 mul add frs add L
          \fi
       0 \psk@fontscale 0.9 mul frs 2 mul add neg rlineto closepath
        gsave
         \pst@usecolor\psfillcolor \tx@setTransparency fill
        grestore
         \ifx\pslinestyle\@none\else
           \pst@number\pslinewidth SLW \pst@usecolor\pslinecolor \tx@setStrokeTransparency \@nameuse{psls@\pslinestyle} stroke
         \fi
       \ifPst@originV 0 0 \else
          xM xShift add \pst@number\psxunit mul
          yM yShift add \pst@number\psyunit mul
       \fi
     xorytext
   \fi
  \fi
\fi}%
\ifPst@markZeros\psdot(#5)\fi
\endgroup\use@pscode\ignorespaces
}%
%
\catcode`\@=\PstAtCode\relax
%
%% END: pst-func.tex
\endinput
%

